% page 142 of the facsimile (image p-141 of the scan)
% block 157.5 x 254.9 mm ; left margin 8.0 mm
\pgc{-12.700mm}{0.000mm}{
\l{\cel{7}132.}
\l{}
\l{}
\l{\cel{1}\sou{emberizo}.- (\sou{zool.}). Genero de uceli, di qua la speco tipa}
\l{\cel{4}esas pasero \textquotedbl{}kornirostra\textquotedbl{}, analoga a la verdrono. - L.}
\l{\cel{4}\sou{emberiza}.}
\l{}
\l{\cel{1}\sou{emblemo}. I. Signo videbla, selektita konvencione por repre-}
\l{\cel{3}esentar ideo, kozo abstraktita, ed ulakaze akompanata da}
\l{\cel{4}subskriburo explika. -II. Figuro qua memorigas kozo abs-}
\l{\cel{4}traktita per ula aludo. - III. Atributo di kozo, uzata}
\l{\cel{4}por reprezentar olca. - DEFIRS.}
\l{}
\l{\cel{1}embolo. (patol.)\cel{2}Korpo stranjera qua produktas la fenomeno}
\l{\cel{4}\textquotedbl{}embolio\textquotedbl{}.}
\l{}
\l{\cel{1}\sou{embolio}. (\sou{patol.}) I. Sanga koagulajo fibrinoza, kuntirata da}
\l{\cel{4}la sango-cirkulado aden le arterii plu mikra,quin lu povas}
\l{\cel{4}obstruktar. - II. Sango koagulajo, naskinta en la veini,}
\l{\cel{4}e kuntirata da la sango-cirkulado ad la kordio.}
\l{}
\l{\cel{1}embotar. (trans.) (tekn.)\cel{2}I.Fasonar, per martelago, metalo}
\l{\cel{4}por formacar sur olca la reliefo de puncuro ube la repro-}
\l{\cel{4}duktendo esas kava. - II. Fasonar, per martelago o saliig-}
\l{\cel{4}ilo, metalo-plako, fero-tolo, por rondigar lu a baseno,}
\l{\cel{4}kasrolo, por igar lu ornivo arkitekturala. - FS.}
\l{}
\l{\cel{1}embracar. ( trans.) Prenar e klemar en sua brakii. - (alud-}
\l{\cel{4}ante kozo).Cirkondar per sua korpo - quaze kovrar.- EFIS.}
\l{}
\l{\cel{1}\sou{embragar}. (\sou{trans.}) Igar komunikar la motoro di mekanismo}
\l{\cel{4}kun la parti di olca quin onu volas movigar, per glit-}
\l{\cel{4}igar kulise rimeno de libera pulio adsur fixa pulio (sam-}
\l{\cel{4}axo e sama-periferia), dento-roto ad piniono, e c.- F.}
\l{}
\l{\cel{1}\sou{embrazuro}. I. Vakuo en parieto. - II. Vakuo en muro, ube}
\l{\cel{4}ajustesas la klapi di pordo, di fenestro. - III. Vakuo}
\l{\cel{4}en la muro di parapeto, qua recevas la boko di kanono. -}
\l{\cel{4}IV. Aperturo di furnelo tra qua pasas la kolo di retorto.}
\l{\cel{4}- V. Vakuo en la masivo di fornego. - EFR.}
\l{}
\l{\sou{embriogenio}.Formaceso, developeso di la embriono che la ani-}
\l{\cel{4}mali. - DEFIRS.}
\l{}
\l{\cel{1}\sou{embriologio}. Parto di la natur-historio qua traktas la for-}
\l{\cel{4}maceso e la developeso di la embriono.}
\l{}
\l{\cel{1}\sou{embriono}. (\sou{Fiziol.}) Jermo naskanta, produktita da la fekund-}
\l{\cel{4}igo. - II. (\sou{metafore}). Intencajo. - III. Jermo fekundig-}
\l{\cel{4}ita, qua komencas developesar en la semino. - DEFIS.}
\l{}
\l{\cel{1}\sou{embrokar}. (\sou{Medicino)}. (\sou{netrans.}) Fomentar, per varsar lente}
\l{\cel{4}liquido grasa sur korpo-parto malada. - DEFIS.}
\l{}
\l{\cel{1}\sou{embuskar}. (\sou{trans.)}\cel{2}Postenizar en loko de ube lu povas sur-}
\l{\cel{4}prizar la enemiko (od ulu altra). - EFIS.}
\l{}
\l{\cel{1}emendar.\cel{1}\sou{(trans}.) I. Plubonigar per modifikar to quo esas}
\l{\cel{4}mala, defektoza. II. (Yuro) Reformar verdikto.-DEFIS.}
}
